Classifier guidance was introduced by \citet{NEURIPS2021_49ad23d1} as a mechanism for steering the
reverse diffusion trajectory toward samples that satisfy a desired condition $y$.  Instead of
modifying the DDPM mean parameter, the formulation is more naturally expressed in terms of the
noise-prediction parameterization
\[
\nabla_{x_t} \log p(x_t)
=
-\frac{1}{\sqrt{1-\bar\alpha_t}}\,\epsilon_\theta(x_t,t),
\]
which relates the score to the model’s predicted noise. 
The authors observed that the classifier gradient
\[
\nabla_{x_t} \log p_\phi(y \mid x_t)
\]
can be used as an additional force that biases the sampling process toward regions where the 
classifier assigns higher likelihood to the target condition.

Using the identity
\[
\nabla_{x_t}\log p(y \mid x_t)
=
\nabla_{x_t}\log p(x_t \mid y)
-
\nabla_{x_t}\log p(x_t),
\]
we note that the term $\nabla_{x_t}\log p(y)$ vanishes because $p(y)$ does not depend on the image 
variable $x_t$.  
This follows directly from Bayes’ rule,
\[
p(y \mid x_t) = \frac{p(x_t \mid y)\,p(y)}{p(x_t)},
\]
whose logarithm yields 
\[
\log p(y \mid x_t)
=
\log p(x_t \mid y)
+
\log p(y)
-
\log p(x_t),
\]
and differentiation with respect to $x_t$ eliminates the constant $\log p(y)$.
Thus the classifier gradient can be written purely as a difference of conditional and 
unconditional data scores.

Substituting the score--noise relation, the guided noise predictor takes the form \citep{NEURIPS2021_49ad23d1, weng2021diffusion}:
\begin{equation}
\tilde{\epsilon}_\theta(x_t,t,y)
=
\epsilon_\theta(x_t,t,y)
+
w\bigl(\epsilon_\theta(x_t,t,y) - \epsilon_\theta(x_t,t)\bigr)
=
(1+w)\,\epsilon_\theta(x_t,t,y)
-
w\,\epsilon_\theta(x_t,t).
\label{eq:guided_eps_classifier}
\end{equation}
Here $w \ge 0$ is the guidance weight.  
Larger values of $w$ increase adherence to the target condition at the cost of reduced sample
diversity.  
This expression makes clear that classifier guidance amplifies the difference between the
conditional and unconditional denoising directions.

\paragraph{Classifier-Free Guidance.}

Classifier-free guidance (CFG) \citep{ho2022classifier} achieves the same effect without requiring a
separate classifier.  
A single diffusion model is trained so that it learns both the conditional and unconditional noise
predictors:
\[
\epsilon_{\mathrm{cond}}(x_t,t,y), \qquad
\epsilon_{\mathrm{uncond}}(x_t,t),
\]
by randomly dropping the conditioning variable $y$ during training.  
With these two predictors, \citet{ho2022classifier} showed that the same guided update used in
classifier guidance can be reproduced purely from the model itself:
\begin{equation}
\tilde{\epsilon}_\theta(x_t,t,y)
=
(1+w)\,\epsilon_{\mathrm{cond}}(x_t,t,y)
-
w\,\epsilon_{\mathrm{uncond}}(x_t,t),
\label{eq:guided_eps_cfg}
\end{equation}
which is algebraically equivalent to applying classifier guidance with weight $w+1$ to an
unconditional diffusion model.  
The hyperparameter $w$ again determines the strength of conditioning: $w=0$ yields standard
conditional generation, while $w>1$ produces extrapolative guidance that strengthens the conditional
signal at the cost of diversity.

This formulation has become the dominant approach in conditional diffusion models, as it retains the
effectiveness of classifier guidance while avoiding the need for a separate classifier and enabling
conditional and unconditional behavior to be learned jointly within one network.